\chapter[15]{Differential Harnack Estimates of LYH-type}
\label{ch:chow-ii-differential-harnack-estimates}

Let \((\mathcal M^n,g(t))\) be a Ricci flow.  Indices are raised using
\(g(t)\), and \(R_{ijkl}\), \(R_{ij}\), and \(R\) denote its Riemann,
Ricci, and scalar curvatures.

\section{The Harnack expression and Ricci solitons}

\begin{definition}[The tensors \(P\) and \(M\)]
\label{def:chow-ii-harnack-p-and-m-tensors}
\source{chow-et-al-ricci-flow-part-ii:page-0287,chow-et-al-ricci-flow-part-ii:page-0288}
\dcref{ch15:eq-15.5,eq-15.9}
\group{Harnack Quadratic}

For a Ricci flow at time \(t>0\), define
\[
P_{kij}=\nabla_kR_{ij}-\nabla_iR_{kj}
\]
and
\[
M_{ij}=\Delta R_{ij}-\frac12\nabla_i\nabla_jR
 +2R^{\ell}{}_{ijp}R_{\ell}{}^p-R_i{}^pR_{pj}+\frac1{2t}R_{ij}.
\]
Then \(P_{kij}=-P_{ikj}\) and
\(P_{ijk}+P_{jki}+P_{kij}=0\).
\end{definition}
\begin{proof}
\sketch The skew symmetry is immediate from the definition.  The cyclic identity is
the sum of the three displayed differences; every covariant derivative of the
Ricci tensor occurs once with each sign.
\end{proof}

\begin{lemma}[Soliton identities for \(P\) and \(M\)]
\label{lem:chow-ii-soliton-harnack-identities}
\source{chow-et-al-ricci-flow-part-ii:page-0287,chow-et-al-ricci-flow-part-ii:page-0288}
\uses{def:chow-ii-harnack-p-and-m-tensors}
\dcref{ch15:eq-15.7,eq-15.10}
\group{Harnack Quadratic}


If a gradient Ricci soliton satisfies
\[
R_{ij}+\dot c(t)g_{ij}=-\nabla_i\nabla_jf,
\]
then
\[
P_{kij}=R^p{}_{kij}\nabla_pf.
\]
For an expanding soliton normalized by \(\dot c(t)=1/(2t)\), one also has
\[
M_{ij}=P_{pji}\nabla^pf.
\]
\end{lemma}
\begin{proof}
\sketch Differentiate the soliton equation and subtract the result with the first and
third differentiating indices interchanged.  Commuting third derivatives of
\(f\) gives the first identity.  Take a divergence of that identity, commute
the two derivatives of the Ricci tensor, use the contracted second Bianchi
identity, and substitute
\(\nabla_i\nabla_jf=-R_{ij}-\dot c,g_{ij}\).  The resulting equality is
\[
0=\Delta R_{ij}-\tfrac12\nabla_i\nabla_jR
 +2R^\ell{}_{ijp}R_\ell{}^p-R_i{}^pR_{pj}
-P_{pji}\nabla^pf+\dot c,R_{ij}.
\]
Putting \(\dot c=1/(2t)\) proves the second assertion.
\end{proof}

\begin{definition}[Hamilton's Harnack quadratic]
\label{def:chow-ii-hamilton-harnack-quadratic}
\source{chow-et-al-ricci-flow-part-ii:page-0289}
\uses{def:chow-ii-harnack-p-and-m-tensors}
\dcref{ch15:eq-15.12}
\group{Harnack Quadratic}


For a two-form \(U\) and one-form \(W\), define
\[
Q(U\oplus W)=M_{ij}W^iW^j+2P_{pij}U^{pi}W^j
 +R_{pijq}U^{pq}U^{ij}.
\]
\end{definition}

\begin{lemma}[The Harnack quadratic is stationary on expanding solitons]
\label{lem:chow-ii-harnack-quadratic-expanding-soliton-nullity}
\source{chow-et-al-ricci-flow-part-ii:page-0289}
\uses{lem:chow-ii-soliton-harnack-identities, def:chow-ii-hamilton-harnack-quadratic}
\dcref{ch15:eq-15.15}
\group{Harnack Quadratic}


On a normalized expanding gradient Ricci soliton,
\[
Q((W\wedge\nabla f)\oplus W)=0
\]
for every \(W\), and the differential of \(Q\) vanishes at
\((W\wedge\nabla f)\oplus W\).
\end{lemma}
\begin{proof}
\sketch The differential in a direction \(A\oplus B\) is
\[
2(M_{ij}-P_{pij}\nabla^pf)W^iB^j
+2(P_{pij}-R_{pijq}\nabla^qf)A^{pi}W^j.
\]
Both coefficients vanish by Lemma~\ref{lem:chow-ii-soliton-harnack-identities},
after using the cyclic identity for \(P\).  Since the value itself vanishes by
the same identities, both conclusions follow.
\end{proof}

\section{Hamilton's matrix estimate and its consequences}

\begin{theorem}[Hamilton's matrix Harnack estimate]
\label{thm:chow-ii-hamilton-matrix-harnack}
\source{chow-et-al-ricci-flow-part-ii:page-0290,chow-et-al-ricci-flow-part-ii:page-0306,chow-et-al-ricci-flow-part-ii:page-0307,chow-et-al-ricci-flow-part-ii:page-0308,chow-et-al-ricci-flow-part-ii:page-0309,chow-et-al-ricci-flow-part-ii:page-0310,chow-et-al-ricci-flow-part-ii:page-0311,chow-et-al-ricci-flow-part-ii:page-0312,chow-et-al-ricci-flow-part-ii:page-0313,chow-et-al-ricci-flow-part-ii:page-0314}
\uses{def:chow-ii-hamilton-harnack-quadratic, prop:chow-ii-bilinear-form-maximum-principle, thm:chow-ii-harnack-quadratic-evolution, lem:chow-ii-harnack-exhaustion-function, lem:chow-ii-harnack-barrier-functions, lem:chow-ii-perturbed-harnack-evolution}
\dcref{ch15:15.1}
\group{Matrix Harnack Estimates}


If \((\mathcal M^n,g(t))\), \(t\in[0,T)\), is complete and has bounded
nonnegative curvature operator, then
\[
Q(U\oplus W)\ge0
\]
for every two-form \(U\) and one-form \(W\).
\end{theorem}
\begin{proof}
\sketch First suppose \(\mathcal M\) is closed and \(\Rm>0\).  For small positive
time the \(R_{ij}/(2t)\) term makes \(Q\) positive.  The evolution formula of
Theorem~\ref{thm:chow-ii-harnack-quadratic-evolution} is a heat inequality
whose reaction term has the null-eigenvector property: when \(Q\ge0\), write
it as a sum of squares
\[
Q(U\oplus W)=\sum_N(Y^N_{ab}U_{ab}+X^N_aW_a)^2.
\]
Substitution into the reaction term gives
\[
\sum_{M,N}(Y^M_{ac}X^N_cW_a-Y^N_{ac}X^M_cW_a
-2Y^M_{ac}Y^N_{bc}U_{ab})^2\ge0.
\]
The bilinear-form maximum principle therefore preserves \(Q\ge0\).

For a complete flow with \(\Rm\ge0\), add \(t^{-1}\phi g\) to \(M\) and
\(\psi I\) to the curvature block.  Lemmas
\ref{lem:chow-ii-harnack-exhaustion-function}--\ref{lem:chow-ii-perturbed-harnack-evolution}
choose positive \(\phi,\psi\), arbitrarily small on any prescribed compact
set, so that the perturbed quadratic is positive near time zero and at spatial
infinity and has positive time derivative at a first null vector.  Such a
first null vector is impossible.  Letting the perturbations tend to zero on
an exhaustion of space-time yields \(Q\ge0\).
\end{proof}

\begin{remark}[No initial-time curvature bound is needed]
\label{rem:chow-ii-harnack-open-initial-time}
\source{chow-et-al-ricci-flow-part-ii:page-0290}
\uses{thm:chow-ii-hamilton-matrix-harnack}
\dcref{ch15:15.2}
\group{Matrix Harnack Estimates}


The conclusion of Theorem~\ref{thm:chow-ii-hamilton-matrix-harnack} remains
valid for a flow on \((0,T)\) whose curvature is bounded on every compact time
interval, even if \(\sup_{\mathcal M}|\Rm(t)|\to\infty\) as \(t\downarrow0\).
\end{remark}
\begin{proof}
\sketch Apply the theorem on \([\varepsilon,T')\), with time translated by
\(t\mapsto t-\varepsilon\).  For fixed \(t>0\), let
\(\varepsilon\downarrow0\); the additional Ricci term is
\(R_{ij}/(2(t-\varepsilon))\), which converges to \(R_{ij}/(2t)\).
Then let \(T'\uparrow T\).
\end{proof}

\begin{corollary}[Trace Harnack estimate and monotonicity of \(tR\)]
\label{cor:chow-ii-trace-harnack-and-tr-monotonicity}
\source{chow-et-al-ricci-flow-part-ii:page-0290,chow-et-al-ricci-flow-part-ii:page-0291}
\uses{thm:chow-ii-hamilton-matrix-harnack}
\dcref{ch15:15.3}
\group{Trace Harnack Estimates}


Under the hypotheses of Theorem~\ref{thm:chow-ii-hamilton-matrix-harnack},
\[
\partial_tR+\frac Rt+2\langle\nabla R,V\rangle+2\Ric(V,V)\ge0.
\]
Consequently \(t\mapsto tR(x,t)\) is nondecreasing and
\(R(x,t_2)\ge(t_1/t_2)R(x,t_1)\) for \(0<t_1\le t_2\).
\end{corollary}
\begin{proof}
\sketch For an orthonormal coframe \(\{\omega^a\}\), sum
\(Q((\omega^a\wedge V)\oplus\omega^a)\ge0\) over \(a\).  The contracted
Bianchi identity and \(\partial_tR=\Delta R+2|\Ric|^2\) give
\[
0\le\tfrac12\partial_tR+\tfrac{R}{2t}-\langle\nabla R,V\rangle+\Ric(V,V).
\]
Replace \(V\) by \(-V\) and multiply by two.  Setting \(V=0\) gives
\(\partial_t(tR)\ge0\), and integration proves the last assertion.
\end{proof}

\begin{corollary}[Trace Harnack estimate for ancient solutions]
\label{cor:chow-ii-ancient-trace-harnack}
\source{chow-et-al-ricci-flow-part-ii:page-0291}
\uses{cor:chow-ii-trace-harnack-and-tr-monotonicity}
\dcref{ch15:15.4}
\group{Trace Harnack Estimates}


If a complete ancient Ricci flow has bounded nonnegative curvature operator,
then
\[
\partial_tR+2\langle\nabla R,V\rangle+2\Ric(V,V)\ge0,
\qquad \partial_tR\ge0.
\]
\end{corollary}
\begin{proof}
\sketch Apply Corollary~\ref{cor:chow-ii-trace-harnack-and-tr-monotonicity} after
translating an arbitrary \(\alpha<t\) to time zero.  Its extra term is
\(R/(t-\alpha)\).  Let \(\alpha\to-\infty\), then take \(V=0\).
\end{proof}

\begin{corollary}[Comparison of scalar curvature at space-time points]
\label{cor:chow-ii-scalar-curvature-spacetime-comparison}
\source{chow-et-al-ricci-flow-part-ii:page-0291,chow-et-al-ricci-flow-part-ii:page-0292}
\uses{cor:chow-ii-trace-harnack-and-tr-monotonicity}
\dcref{ch15:15.5}
\group{Trace Harnack Estimates}


Under the hypotheses of Theorem~\ref{thm:chow-ii-hamilton-matrix-harnack},
\[
\partial_t\log R-\frac12|\nabla\log R|^2+\frac1t\ge0,
\]
and, for \(t_2>t_1>0\),
\[
\frac{R(x_2,t_2)}{R(x_1,t_1)}\ge\frac{t_1}{t_2}
\exp\!\left(-\frac{d_{g(t_1)}(x_1,x_2)^2}{2(t_2-t_1)}\right).
\]
\end{corollary}
\begin{proof}
\sketch Nonnegative curvature operator gives \(0\le\Ric\le Rg\).  Insert
\(V=-\nabla R/(2R)\) into the trace estimate.  For a path
\(\gamma:[t_1,t_2]\to\mathcal M\), the resulting inequality and
\(2\langle X,Y\rangle\le|X|^2+|Y|^2\) imply
\[
\frac d{dt}\log(tR(\gamma(t),t))\ge-\frac12|\dot\gamma(t)|_{g(t)}^2.
\]
Since \(\Ric\ge0\), \(g(t)\le g(t_1)\).  Integrate along a constant-speed
\(g(t_1)\)-minimizing geodesic and minimize its energy.
\end{proof}

\begin{corollary}[Ancient scalar-curvature comparison]
\label{cor:chow-ii-ancient-scalar-curvature-spacetime-comparison}
\source{chow-et-al-ricci-flow-part-ii:page-0292}
\uses{cor:chow-ii-ancient-trace-harnack}
\dcref{ch15:15.6}
\group{Trace Harnack Estimates}


For a complete ancient solution with bounded nonnegative curvature operator,
\[
\frac{R(x_2,t_2)}{R(x_1,t_1)}\ge
\exp\!\left(-\frac{d_{g(t_1)}(x_1,x_2)^2}{2(t_2-t_1)}\right)
\]
whenever \(t_2>t_1\).
\end{corollary}
\begin{proof}
\sketch Repeat the path integration in the preceding proof using the ancient trace
estimate, which has no \(1/t\) term.
\end{proof}

\section{The two-dimensional estimates}

\begin{proposition}[Two-dimensional trace Harnack estimate]
\label{prop:chow-ii-surface-trace-harnack-positive-curvature}
\source{chow-et-al-ricci-flow-part-ii:page-0292,chow-et-al-ricci-flow-part-ii:page-0293}
\dcref{ch15:15.7}
\group{Surface Harnack Estimates}

For a complete surface Ricci flow with bounded positive scalar curvature,
\[
\partial_t\log R-|\nabla\log R|^2+\frac1t
=\Delta\log R+R+\frac1t>0.
\]
Equivalently, for every vector field \(V\),
\[
\partial_tR+\frac Rt+2\langle\nabla R,V\rangle+R|V|^2>0.
\]
\end{proposition}
\begin{proof}
\sketch Put \(H=\Delta\log R+R\).  Direct differentiation using
\(\partial_tR=\Delta R+R^2\) yields
\[
\partial_tH\ge\Delta H+2\langle\nabla\log R,\nabla H\rangle+H^2.
\]
The maximum principle and the solution \(-1/t\) of \(q'=q^2\) give
\(H+1/t>0\).  Completing the square in \(V\) proves the equivalent form.
\end{proof}

\begin{remark}[Weak surface trace estimate]
\label{rem:chow-ii-surface-trace-harnack-nonnegative-curvature}
\source{chow-et-al-ricci-flow-part-ii:page-0293}
\uses{prop:chow-ii-surface-trace-harnack-positive-curvature}
\dcref{ch15:15.8}
\group{Surface Harnack Estimates}


For nonnegative scalar curvature, the quadratic trace Harnack estimate remains
valid with weak inequality.
\end{remark}
\begin{proof}
\sketch Apply the positive-curvature estimate to \(R+\varepsilon\) in the standard
strictly positive barrier argument and let \(\varepsilon\downarrow0\).
\end{proof}

\begin{corollary}[Surface monotonicity of \(tR\)]
\label{cor:chow-ii-surface-tr-monotonicity}
\source{chow-et-al-ricci-flow-part-ii:page-0293}
\uses{rem:chow-ii-surface-trace-harnack-nonnegative-curvature}
\dcref{ch15:15.9}
\group{Surface Harnack Estimates}


On a surface Ricci flow with \(R\ge0\), \(tR(x,t)\) is nondecreasing.  For an
ancient flow, \(R(x,t)\) is nondecreasing.
\end{corollary}
\begin{proof}
\sketch Set \(V=0\) in the weak trace estimate.  For an ancient flow translate the
initial time to \(\alpha\) and let \(\alpha\to-\infty\).
\end{proof}

\begin{proposition}[Two-dimensional matrix Harnack estimate]
\label{prop:chow-ii-surface-matrix-harnack}
\source{chow-et-al-ricci-flow-part-ii:page-0293}
\dcref{ch15:15.10}
\group{Surface Harnack Estimates}

For a complete surface Ricci flow on \([0,T)\) with bounded positive scalar
curvature,
\[
\nabla_i\nabla_j\log R+\frac12Rg_{ij}+\frac1{2t}g_{ij}>0.
\]
\end{proposition}
\begin{proof}
\sketch Let \(L=\log R\) and
\(H_{ij}=\nabla_i\nabla_jL+\tfrac12(R+t^{-1})g_{ij}\).  Calculation gives
\[
\partial_tH_{ij}=\Delta H_{ij}+2\nabla_kH_{ij}\nabla_kL
+2H_{ik}H_{jk}-(2t^{-1}+3R)H_{ij}+R(\operatorname{tr}H)g_{ij}.
\]
The tensor is positive for small time.  At a first null eigenvector every
reaction term surviving contraction with that vector is nonnegative, so the
tensor maximum principle precludes such a first zero.
\end{proof}

\begin{remark}[Trace, weak, and equality cases on surfaces]
\label{rem:chow-ii-surface-matrix-harnack-consequences}
\source{chow-et-al-ricci-flow-part-ii:page-0293}
\uses{prop:chow-ii-surface-matrix-harnack, prop:chow-ii-surface-trace-harnack-positive-curvature}
\dcref{ch15:15.11}
\group{Surface Harnack Estimates}


Tracing the surface matrix estimate gives the trace estimate.  A flow defined
only on \((0,T)\) satisfies the weak matrix inequality, and equality can occur
for positively curved expanding gradient solitons limiting to a cone at time
zero.
\end{remark}

\begin{remark}[Evolution formula for the surface matrix]
\label{rem:chow-ii-surface-matrix-evolution-reference}
\source{chow-et-al-ricci-flow-part-ii:page-0293}
\dcref{ch15:15.12}
\group{Surface Harnack Estimates}

The evolution formula used in Proposition~\ref{prop:chow-ii-surface-matrix-harnack}
is the tensorial mechanism behind the two-dimensional matrix estimate.
\end{remark}

\begin{proposition}[Equivalence of the two surface matrix estimates]
\label{prop:chow-ii-equivalent-surface-matrix-harnack-forms}
\source{chow-et-al-ricci-flow-part-ii:page-0294,chow-et-al-ricci-flow-part-ii:page-0295}
\uses{def:chow-ii-hamilton-harnack-quadratic, prop:chow-ii-surface-matrix-harnack}
\dcref{ch15:15.13}
\group{Surface Harnack Estimates}


In dimension two with \(R>0\), Hamilton's inequality \(Q\ge0\) is equivalent
to
\[
\nabla_i\nabla_jR+\left(\frac{R^2}{2}+\frac{R}{2t}\right)g_{ij}
-\frac1R\nabla_iR\nabla_jR\ge0,
\]
and hence to the matrix inequality for \(\log R\).
\end{proposition}
\begin{proof}
\sketch Use \(R_{ij}=Rg_{ij}/2\) and write every two-form as
\(U=f\,e^1\wedge e^2\).  Completing the square in \(f\) gives
\[
2Q=\nabla_V\nabla_VR+\left(\frac{R^2}{2}+\frac{R}{2t}\right)|V|^2
-\frac{\langle\nabla R,V\rangle^2}{R}
+\frac{(\langle\nabla R,V\rangle-2fR)^2}{R}.
\]
Minimizing in \(f\), then varying \(V\), yields the displayed tensor
inequality.  Dividing by \(R\) is precisely the Hessian identity for
\(\log R\).
\end{proof}

\section{Evolution and maximum-principle proof}

\begin{proposition}[Maximum principle for bilinear forms]
\label{prop:chow-ii-bilinear-form-maximum-principle}
\source{chow-et-al-ricci-flow-part-ii:page-0295,chow-et-al-ricci-flow-part-ii:page-0296,chow-et-al-ricci-flow-part-ii:page-0297,chow-et-al-ricci-flow-part-ii:page-0298}
\dcref{ch15:15.14}
\group{Tensor Maximum Principles}

Let \(\mathcal V\) be a finite-rank metric vector bundle over a closed
time-dependent Riemannian manifold.  Suppose \(F_t\) is fiber preserving,
locally Lipschitz on bounded sets, and has the null-eigenvector property
\[
S\ge0,\quad S(v,v)=0\quad\Longrightarrow\quad F_t(S)(v,v)\ge0.
\]
Suppose a symmetric bilinear form \(S(t)\) satisfies, after extending every
fiber vector \(v\) near a chosen point,
\[
\partial_tS(v,v)\ge\Delta S(v,v)+F_t(S)(v,v),
\]
while \(\partial_th\ge-C_1h\), \(|\nabla v|\le C_2|v|\), and
\(|(\partial_t-\Delta)v|\le C_3|v|\).  If \(S(0)\ge0\), then
\(S(t)\ge0\) for all subsequent times.
\end{proposition}
\begin{proof}
\sketch On a short interval beginning at \(t_0\), put
\(S_\varepsilon=S+\varepsilon(\delta+t-t_0)h\).  Local Lipschitz continuity
and the three extension bounds give, for sufficiently small \(\delta>0\),
\[
\partial_tS_\varepsilon(v,v)\ge\Delta S_\varepsilon(v,v)
+F_t(S_\varepsilon)(v,v)+\tfrac12\varepsilon|v|^2.
\]
If \(S_\varepsilon\) first acquired a null vector, the spatial second
derivative test, the time first derivative test, and the null-eigenvector
property would contradict this strict inequality.  Thus
\(S_\varepsilon>0\) on the interval.  Let \(\varepsilon\downarrow0\), then
iterate the uniform short interval finitely many times.
\end{proof}

\begin{remark}[Comparison with the tensor maximum principle]
\label{rem:chow-ii-bilinear-form-maximum-principle-comparison}
\source{chow-et-al-ricci-flow-part-ii:page-0296}
\dcref{ch15:15.15}
\group{Tensor Maximum Principles}

Proposition~\ref{prop:chow-ii-bilinear-form-maximum-principle} is the vector
bundle version of the weak maximum principle for symmetric tensors, with
explicit control of the chosen local vector extensions.
\end{remark}

\begin{theorem}[Evolution of Hamilton's matrix Harnack quadratic]
\label{thm:chow-ii-harnack-quadratic-evolution}
\source{chow-et-al-ricci-flow-part-ii:page-0299,chow-et-al-ricci-flow-part-ii:page-0300,chow-et-al-ricci-flow-part-ii:page-0301,chow-et-al-ricci-flow-part-ii:page-0302,chow-et-al-ricci-flow-part-ii:page-0303,chow-et-al-ricci-flow-part-ii:page-0304,chow-et-al-ricci-flow-part-ii:page-0305,chow-et-al-ricci-flow-part-ii:page-0306}
\uses{def:chow-ii-hamilton-harnack-quadratic, lem:chow-ii-evolution-of-p-tensor, lem:chow-ii-evolution-of-m-tensor}
\dcref{ch15:15.16}
\group{Harnack Evolution Equations}


At a point, extend \(U\) and \(W\) so that
\[
(\nabla_t-\Delta)W_a=t^{-1}W_a,\quad
(\nabla_t-\Delta)U_{ab}=0,\quad \nabla_aW_b=0,
\]
\[
\nabla_aU_{bc}=\tfrac12(R_{ab}W_c-R_{ac}W_b)
+\tfrac1{4t}(g_{ab}W_c-g_{ac}W_b).
\]
Then
\[
\begin{aligned}
(\nabla_t-\Delta)Q={}&2R_{acdb}M_{cd}W_aW_b-2P_{acd}P_{bdc}W_aW_b\\
&+8R_{adec}P_{dbe}U_{ab}W_c+4R_{aefc}R_{befd}U_{ab}U_{cd}\\
&+(P_{abc}W_c+R_{abdc}U_{cd})
(P_{abe}W_e+R_{abfe}U_{ef}).
\end{aligned}
\]
\end{theorem}
\begin{proof}
\sketch Apply \(\nabla_t-\Delta\) to the three summands defining \(Q\), use the
product rule, and substitute the evolution equations for \(R_{abcd}\),
\(P_{abc}\), and \(M_{ab}\).  The prescribed derivatives of \(U,W\) cancel
all first-order terms.  The remaining terms linear in \(M,P,R\) are the first
two lines displayed above.  The residual terms are
\[
P_{abc}P_{abe}W_cW_e+2P_{abc}R_{abfe}W_cU_{ef}
+R_{abdc}R_{abfe}U_{cd}U_{ef},
\]
which factor as the last line.
\end{proof}

\begin{remark}[Soliton origin of the extension equations]
\label{rem:chow-ii-harnack-extension-soliton-motivation}
\source{chow-et-al-ricci-flow-part-ii:page-0299,chow-et-al-ricci-flow-part-ii:page-0300}
\dcref{ch15:15.17}
\group{Harnack Evolution Equations}

The spatial extension equation for \(U\) is obtained by taking
\(U=W\wedge\nabla f\) on a normalized expanding gradient soliton and using
\(\nabla^2f=-\Ric-g/(2t)\).  The time equations are the corresponding
parabolic extensions at the chosen point.
\end{remark}

\begin{lemma}[Evolution of the \(P\)-tensor]
\label{lem:chow-ii-evolution-of-p-tensor}
\source{chow-et-al-ricci-flow-part-ii:page-0300,chow-et-al-ricci-flow-part-ii:page-0301}
\uses{def:chow-ii-harnack-p-and-m-tensors}
\dcref{ch15:15.18}
\group{Harnack Evolution Equations}


\[
(\nabla_t-\Delta)P_{abc}=2R_{adeb}P_{dec}+2R_{adec}P_{dbe}
+2R_{bdec}P_{ade}-2R_{de}\nabla_dR_{abec}.
\]
\end{lemma}
\begin{proof}
\sketch Apply the heat operator to
\(P_{abc}=\nabla_aR_{bc}-\nabla_bR_{ac}\).  Commute it through each
covariant derivative, substitute
\((\nabla_t-\Delta)R_{bc}=2R_{bdec}R_{de}\), and use the second Bianchi
identity.  Grouping the three Ricci-derivative differences gives the three
\(P\)-terms, while the differentiated curvature term is
\(-2R_{de}\nabla_dR_{abec}\).
\end{proof}

\begin{lemma}[Evolution of the \(M\)-tensor]
\label{lem:chow-ii-evolution-of-m-tensor}
\source{chow-et-al-ricci-flow-part-ii:page-0301,chow-et-al-ricci-flow-part-ii:page-0302,chow-et-al-ricci-flow-part-ii:page-0303,chow-et-al-ricci-flow-part-ii:page-0304}
\uses{def:chow-ii-harnack-p-and-m-tensors, lem:chow-ii-evolution-of-p-tensor}
\dcref{ch15:15.19}
\group{Harnack Evolution Equations}


\[
\begin{aligned}
(\nabla_t-\Delta)M_{ab}={}&2R_{acdb}M_{cd}
+2R_{cd}(\nabla_cP_{dab}+\nabla_cP_{dba})\\
&+P_{cda}P_{cdb}-2P_{acd}P_{bdc}
+2R_{cd}R_{ce}R_{adeb}-\frac1{2t^2}R_{ab}.
\end{aligned}
\]
\end{lemma}
\begin{proof}
\sketch Contracted Bianchi and one commutation give
\[
M_{ab}=\nabla_cP_{cab}+R_{acdb}R_{cd}+\frac1{2t}R_{ab}.
\]
Apply \(\nabla_t-\Delta\).  Use Lemma~\ref{lem:chow-ii-evolution-of-p-tensor},
the curvature and Ricci evolution equations, and commute the derivative in
\(\nabla_cP_{cab}\).  The derivative terms combine into the displayed
\(2R_{cd}(\nabla_cP_{dab}+\nabla_cP_{dba})\); the quadratic derivative terms
combine into \(P_{cda}P_{cdb}-2P_{acd}P_{bdc}\); the algebraic curvature
terms combine into the first and penultimate terms.  Differentiating
\(1/(2t)\) supplies the final term.
\end{proof}

\begin{remark}[Equivalent form of the \(P^2\)-terms]
\label{rem:chow-ii-m-evolution-equivalent-p-terms}
\source{chow-et-al-ricci-flow-part-ii:page-0301}
\dcref{ch15:15.20}
\group{Harnack Evolution Equations}

The cyclic identity for \(P\) gives
\[
P_{cda}P_{cdb}-2P_{acd}P_{bdc}
=2P_{acd}P_{bcd}-4P_{acd}P_{bdc}.
\]
\end{remark}

\begin{example}[Unconstrained evolution of the Harnack quadratic]
\label{ex:chow-ii-unconstrained-harnack-quadratic-evolution}
\source{chow-et-al-ricci-flow-part-ii:page-0306}
\uses{thm:chow-ii-harnack-quadratic-evolution}
\dcref{ch15:15.21}
\group{Harnack Evolution Equations}


For arbitrary space-time extensions of \(U,W\), the evolution of \(Q\) is the
formula of Theorem~\ref{thm:chow-ii-harnack-quadratic-evolution} plus the
product-rule terms obtained from the differences between their actual
derivatives and the four prescribed derivatives in that theorem.
\end{example}
\begin{proof}
\sketch Apply the product rule to each occurrence of \(U\) and \(W\).  Subtract and
add the four prescribed derivatives.  The added terms reproduce the theorem;
linearity of first derivatives and of \(\nabla_t-\Delta\) leaves exactly the
stated correction terms.
\end{proof}

\begin{lemma}[Distance-like exhaustion with uniform derivatives]
\label{lem:chow-ii-harnack-exhaustion-function}
\source{chow-et-al-ricci-flow-part-ii:page-0309}
\dcref{ch15:15.22}
\group{Noncompact Harnack Argument}

If a complete noncompact Ricci flow on \([0,T]\) satisfies
\(|\Rm|+|\nabla\Ric|\le k_0\), then for each \(O\in\mathcal M\) there is a
smooth \(h\) and \(C=C(n,k_0,T)\) such that
\[
d_{g(t)}(O,x)+1\le h(x)\le C(d_{g(t)}(O,x)+1),
\quad |\nabla h|+|\nabla^2h|\le C.
\]
\end{lemma}
\begin{proof}
\sketch Choose a distance-like function with bounded first and second derivatives for
\(g(0)\).  The curvature bound makes all \(g(t)\) uniformly equivalent on
\([0,T]\).  Moreover \(\partial_t\Gamma=\nabla\Ric* g^{-1}\), so the
connection bound follows by integration.  Rescaling \(h\) absorbs the uniform
constants and proves all estimates.
\end{proof}

\begin{lemma}[Barrier functions for the perturbed quadratic]
\label{lem:chow-ii-harnack-barrier-functions}
\source{chow-et-al-ricci-flow-part-ii:page-0310}
\uses{lem:chow-ii-harnack-exhaustion-function}
\dcref{ch15:15.23}
\group{Noncompact Harnack Argument}


Given \(C,\eta>0\) and compact \(\Omega\subset\mathcal M\), there are smooth
positive \(\phi(x,t)\), \(\psi(t)\) such that \(\phi,\psi\le\eta\) on
\(\Omega\times[0,T]\), \(\phi\to\infty\) at spatial infinity, and
\[
(\nabla_t-\Delta)\phi>C\phi,\qquad \partial_t\psi>C\psi,
\qquad \phi\ge C\psi.
\]
\end{lemma}
\begin{proof}
\sketch With \(h\) from Lemma~\ref{lem:chow-ii-harnack-exhaustion-function}, take
\(\phi=\varepsilon e^{At}h\) and \(\psi=\delta e^{Bt}\).  Choose
\(A,B>C\), then choose \(\varepsilon\) small enough on \(\Omega\), and
finally \(\delta\) small enough that \(\psi\le\eta\) and
\(C\psi\le\phi\).  The derivative inequalities follow from
\(h\ge1\) and \(\Delta h\le C\).
\end{proof}

\begin{lemma}[Evolution inequality for the perturbed Harnack quadratic]
\label{lem:chow-ii-perturbed-harnack-evolution}
\source{chow-et-al-ricci-flow-part-ii:page-0310,chow-et-al-ricci-flow-part-ii:page-0311,chow-et-al-ricci-flow-part-ii:page-0312,chow-et-al-ricci-flow-part-ii:page-0313,chow-et-al-ricci-flow-part-ii:page-0314}
\uses{thm:chow-ii-harnack-quadratic-evolution, lem:chow-ii-harnack-barrier-functions}
\dcref{ch15:15.24}
\group{Noncompact Harnack Argument}


Define
\[
\widehat{M}=M+t^{-1}\phi g,\qquad \widehat{\Rm}=\Rm+\psi I,
\]
and let \(\widehat Q\) be obtained from \(Q\) by these replacements.  If
\(U,W\) have the extensions in
Theorem~\ref{thm:chow-ii-harnack-quadratic-evolution}, then
\((\nabla_t-\Delta)\widehat Q\) is bounded below by that theorem's reaction
expression with \(M,\Rm\) replaced by \(\widehat{M},\widehat{\Rm}\), plus
\[
\frac1t\left((\nabla_t-\Delta)\phi+\frac\phi t-\frac C t\psi-C\phi\right)|W|^2
+(\partial_t\psi-C\psi)|U|^2.
\]
\end{lemma}
\begin{proof}
\sketch Write \(\widehat Q=Q+t^{-1}\phi|W|^2+\psi|U|^2\) and apply the heat
operator.  The prescribed extensions give
\(|\nabla U|^2\le Ct^{-2}|W|^2\).  Substitute the evolution theorem into the
\(Q\)-term and replace \(M,\Rm\) by their hatted values.  Every replacement
error is bounded by \(Ct^{-1}\phi|W|^2+Ct^{-2}\psi|W|^2+C\psi|U|^2\),
which yields the displayed lower bound.
\end{proof}

\section{The minimizing two-tensor formulation}

Assume in this section that the curvature operator is positive, and denote its
inverse on two-forms by \(S_{abcd}\).  Completing the square in \(U\) gives
\[
Q(U\oplus W)=Z_{ab}W_aW_b
+R_{abcd}(U_{ab}-S_{abij}P_{ijp}W_p)
 (U_{dc}-S_{dck\ell}P_{k\ell q}W_q),
\]
where
\[
Z_{ab}=M_{ab}-S_{ijk\ell}P_{ija}P_{\ell kb}.
\]
Thus \(Q\ge0\) is equivalent to \(Z\ge0\).

\begin{theorem}[Evolution of the minimizing Harnack two-tensor]
\label{thm:chow-ii-minimizing-harnack-two-tensor-evolution}
\source{chow-et-al-ricci-flow-part-ii:page-0315,chow-et-al-ricci-flow-part-ii:page-0316}
\uses{def:chow-ii-hamilton-harnack-quadratic, lem:chow-ii-evolution-of-p-tensor, lem:chow-ii-evolution-of-m-tensor}
\dcref{ch15:15.25}
\group{Minimizing Harnack Tensor}


As long as \(Z\ge0\), write \(Q\) as a sum of squares with coefficients
\(Y^N_{ab},X^N_a\).  Define
\[
K_{avtu}=\nabla_vP_{tua}-S_{ijrs}P_{ija}\nabla_vR_{rstu}
+R_{tuaw}R_{vw}+\frac1{2t}R_{tuav}
\]
and
\[
L_a^{NM}=2Y_{id}^NY_{jd}^MS_{ijk\ell}P_{k\ell a}
+Y_{ah}^NX_h^M-Y_{ah}^MX_h^N.
\]
Then
\[
(\nabla_t-\Delta)Z_{ab}=2S_{mntu}K_{avut}K_{bvmn}
+\sum_{N,M}L_a^{NM}L_b^{NM}-\frac2tZ_{ab}.
\]
\end{theorem}
\begin{proof}
\sketch Differentiate \(Z=M-S(P,P)\).  The inverse identity \(S\Rm=I\) gives
\((\nabla_t-\Delta)S=-S((\nabla_t-\Delta)\Rm)S\) together with the two
gradient-square correction terms.  Substitute the evolutions of \(M,P,\Rm\)
and complete the derivative terms into \(2S(K,K)\).  The remaining algebraic
terms are expanded using
\(R=\sum_NY^N\otimes Y^N\),
\(P=\sum_NY^N\otimes X^N\), and
\(M=\sum_NX^N\otimes X^N\); they are exactly
\(\sum_{N,M}L^{NM}\otimes L^{NM}-2Z/t\).
\end{proof}

\begin{remark}[Trace of the \(K\)-tensor]
\label{rem:chow-ii-k-tensor-trace-is-z}
\source{chow-et-al-ricci-flow-part-ii:page-0316}
\uses{thm:chow-ii-minimizing-harnack-two-tensor-evolution}
\dcref{ch15:15.26}
\group{Minimizing Harnack Tensor}


The tensor \(K\) is skew in its last two indices and satisfies
\[
g^{vt}K_{avtu}=Z_{ua}.
\]
Moreover \(S_{mntu}K_{avut}K_{bvmn}\) is nonnegative as a symmetric
two-tensor.
\end{remark}
\begin{proof}
\sketch Contract the definition of \(K\), use
\(\nabla_vR_{rsvu}=P_{urs}\), and use the divergence formula for \(P\) in
the proof of Lemma~\ref{lem:chow-ii-evolution-of-m-tensor}.  Positivity follows
by writing the positive operator \(S\) as a sum of squares on two-forms.
\end{proof}

\begin{proposition}[General evolution of the minimizing Harnack tensor]
\label{prop:chow-ii-general-z-tensor-evolution}
\source{chow-et-al-ricci-flow-part-ii:page-0317}
\uses{thm:chow-ii-minimizing-harnack-two-tensor-evolution}
\dcref{ch15:15.27}
\group{Minimizing Harnack Tensor}


Without assuming \(Z\ge0\),
\[
\begin{aligned}
(\nabla_t-\Delta)Z_{ab}={}&2S_{mntu}K_{avut}K_{bvmn}-\frac2tZ_{ab}
+2R_{cabd}Z_{cd}-2P_{acd}P_{bdc}\\
&+2R_{acbd}S_{ijk\ell}P_{ijc}P_{k\ell d}
-4S_{ijk\ell}R_{idae}P_{k\ell b}P_{dje}\\
&-4S_{ijk\ell}R_{idbe}P_{k\ell a}P_{dje}
+4R_{iepf}R_{jeqf}S_{ijk\ell}P_{k\ell a}S_{pqmn}P_{mnb}.
\end{aligned}
\]
\end{proposition}
\begin{proof}
\sketch Repeat the differentiation of \(M-S(P,P)\) in the preceding theorem without
using a sum-of-squares decomposition of \(Q\).  Completing only the derivative
terms gives \(2S(K,K)\); direct contraction of all remaining curvature and
\(P\)-terms gives the displayed six algebraic terms.
\end{proof}

\begin{lemma}[Sum-of-squares identity for the \(Z\)-reaction term]
\label{lem:chow-ii-z-reaction-sum-of-squares}
\source{chow-et-al-ricci-flow-part-ii:page-0317}
\uses{prop:chow-ii-general-z-tensor-evolution}
\dcref{ch15:15.28}
\group{Minimizing Harnack Tensor}


When \(\Rm>0\) and \(Z\ge0\), the six algebraic terms after \(-2Z/t\) in
Proposition~\ref{prop:chow-ii-general-z-tensor-evolution} equal
\(\sum_{N,M}L_a^{NM}L_b^{NM}\).
\end{lemma}
\begin{proof}
\sketch Use the sum-of-squares coefficients
\(R=\sum_NY^N\otimes Y^N\),
\(P=\sum_NY^N\otimes X^N\), and
\(M=\sum_NX^N\otimes X^N\).  Substitute them into the six terms and expand.
The \(YYSP\)-squares give the terms quadratic in \(S(P,P)\), the \(YX\)-squares
give the \(RZ-P^2\) terms, and their cross terms give the two terms linear in
\(S(P,P)\).  These are precisely the expansion of
\(\sum_{N,M}L^{NM}\otimes L^{NM}\).
\end{proof}

\begin{lemma}[Characterization of the minimizing two-form]
\label{lem:chow-ii-characterization-minimizing-two-form}
\source{chow-et-al-ricci-flow-part-ii:page-0318}
\dcref{ch15:15.29}
\group{Minimizing Harnack Tensor}

For \(\Rm>0\),
\[
P_{abc}W_c+R_{abdc}U_{cd}=0
\quad\Longleftrightarrow\quad
U_{ab}=S_{abij}P_{ijp}W_p.
\]
\end{lemma}
\begin{proof}
\sketch Apply \(S\) to the first equality and use \(S\Rm=I\) on two-forms.  The
converse follows by applying \(\Rm\) to the second equality.
\end{proof}

\begin{remark}[Three-dimensional lower bound for the \(K\)-term]
\label{rem:chow-ii-three-dimensional-k-term-lower-bound}
\source{chow-et-al-ricci-flow-part-ii:page-0319}
\dcref{ch15:15.30}
\group{Minimizing Harnack Tensor}

In dimension three, if
\(A_{ab}=S_{mntu}K_{avtu}K_{bvmn}\), then
\[
A_{ab}\ge(\Ric^{-1})^{cd}Z_{ac}Z_{bd}.
\]
\end{remark}
\begin{proof}
\sketch Diagonalize \(\Ric\) in an orthonormal frame.  Three-dimensional curvature
has no component with three distinct indices, and
\(4S_{ijji}=R_{ijji}^{-1}\).  Expanding \(A(W,W)\), discard nonnegative
terms and apply
\(x^2/a+y^2/b\ge(x+y)^2/(a+b)\).  The trace identity
\(g^{vt}K_{avtu}=Z_{ua}\) identifies the resulting sums with
\((\Ric^{-1})^{cd}Z_{ac}Z_{bd}W_aW_b\).
\end{proof}

\begin{corollary}[Quadratic lower evolution bound for \(Z\) in dimension three]
\label{cor:chow-ii-three-dimensional-z-quadratic-evolution-bound}
\source{chow-et-al-ricci-flow-part-ii:page-0320}
\uses{thm:chow-ii-minimizing-harnack-two-tensor-evolution, rem:chow-ii-three-dimensional-k-term-lower-bound}
\dcref{ch15:15.31}
\group{Minimizing Harnack Tensor}


In dimension three, as long as \(Z\ge0\),
\[
(\nabla_t-\Delta)Z_{ab}\ge
2(\Ric^{-1})^{cd}Z_{ac}Z_{bd}
+\sum_{N,M}L_a^{NM}L_b^{NM}-\frac2tZ_{ab}.
\]
\end{corollary}
\begin{proof}
\sketch Substitute the lower bound of
Remark~\ref{rem:chow-ii-three-dimensional-k-term-lower-bound} into
Theorem~\ref{thm:chow-ii-minimizing-harnack-two-tensor-evolution}.
\end{proof}

\section{Ancient solutions attaining maximum scalar curvature}

For an ancient solution define
\[
\overline M_{ij}=\Delta R_{ij}-\frac12\nabla_i\nabla_jR
+2R^p{}_{\ell ij}R^\ell{}_p-R_i{}^pR_{pj}
\]
and let \(\overline Q\) be Hamilton's quadratic with \(M\) replaced by
\(\overline M\).

\begin{theorem}[An ancient maximizer is a steady gradient soliton]
\label{thm:chow-ii-ancient-positive-curvature-maximizer-steady-soliton}
\source{chow-et-al-ricci-flow-part-ii:page-0320,chow-et-al-ricci-flow-part-ii:page-0321,chow-et-al-ricci-flow-part-ii:page-0322,chow-et-al-ricci-flow-part-ii:page-0323,chow-et-al-ricci-flow-part-ii:page-0324,chow-et-al-ricci-flow-part-ii:page-0325,chow-et-al-ricci-flow-part-ii:page-0326,chow-et-al-ricci-flow-part-ii:page-0327}
\uses{cor:chow-ii-ancient-matrix-harnack, lem:chow-ii-ancient-harnack-nullspace-propagation, lem:chow-ii-nullspace-map-is-wedge-map, lem:chow-ii-evolution-of-p-tensor, lem:chow-ii-evolution-of-m-tensor}
\dcref{ch15:15.33}
\group{Ancient Solutions and Solitons}


Let a simply connected manifold carry a complete ancient Ricci flow with
strictly positive curvature operator and curvature uniformly bounded in
space-time.  If \(R\) attains its space-time maximum, then the flow is a
steady gradient Ricci soliton.
\end{theorem}
\begin{proof}
\sketch At a maximum of \(R\),
\(0=\partial_tR=2\sum_a\overline Q(0\oplus e_a)\), so the null space of
\(\overline Q\) has dimension \(n\).  Lemma
\ref{lem:chow-ii-ancient-harnack-nullspace-propagation} makes this true at all
earlier points and times.  Positivity of \(\Rm\) makes projection of this null
space to one-forms injective, so it is the graph of a linear map \(L\).
Lemma~\ref{lem:chow-ii-nullspace-map-is-wedge-map} gives a vector field \(V\)
with \(L(W)=V\wedge W\).

First variation of \(\overline Q\) in its null directions gives
\[
P_{ijk}+R_{ijk\ell}V^\ell=0,
\qquad \overline M_{ik}-V^jP_{ijk}=0.
\]
Taking a divergence yields
\(R_{ijk\ell}(\nabla_jV_\ell-R_{j\ell})=0\).  Its skew part and positivity
of the curvature operator show that \(V\) is closed, hence \(V=\nabla f\) by
simple connectivity.  Applying the heat operator to the second first-variation
identity and using the evolutions of \(M\) and \(P\) yields
\(R_{ijk\ell}T_{im}T_{m\ell}=0\), where
\(T=\Ric-\nabla V\).  Diagonalizing \(T\) and using \(\Rm>0\) forces every
eigenvalue of \(T\) to vanish.  Hence \(\Ric=\nabla^2f\); after replacing
\(f\) by \(-f\) according to the flow convention, this is the steady gradient
soliton equation.
\end{proof}

\begin{remark}[Why shrinking solitons do not satisfy the maximum hypothesis]
\label{rem:chow-ii-shrinking-soliton-no-spacetime-r-maximum}
\source{chow-et-al-ricci-flow-part-ii:page-0320}
\dcref{ch15:15.34}
\group{Ancient Solutions and Solitons}

A shrinking sphere, and more generally a shrinking soliton, has scalar
curvature tending to infinity at its singular time, so its space-time scalar
curvature supremum is not attained.
\end{remark}

\begin{corollary}[Matrix Harnack estimate for ancient solutions]
\label{cor:chow-ii-ancient-matrix-harnack}
\source{chow-et-al-ricci-flow-part-ii:page-0321}
\uses{thm:chow-ii-hamilton-matrix-harnack}
\dcref{ch15:15.35}
\group{Ancient Solutions and Solitons}


If a complete ancient Ricci flow has bounded nonnegative curvature operator,
then \(\overline Q(U\oplus W)\ge0\) for every \(U,W\).
\end{corollary}
\begin{proof}
\sketch Translate an arbitrary ancient time \(\alpha\) to zero and apply
Theorem~\ref{thm:chow-ii-hamilton-matrix-harnack}.  The only
\(\alpha\)-dependent term is \(R_{ij}/(2(t-\alpha))\).  Let
\(\alpha\to-\infty\).
\end{proof}

\begin{remark}[An alternative proof of the ancient rigidity theorem]
\label{rem:chow-ii-alternative-ancient-rigidity-proof}
\source{chow-et-al-ricci-flow-part-ii:page-0321}
\dcref{ch15:15.36}
\group{Ancient Solutions and Solitons}

Theorem~\ref{thm:chow-ii-ancient-positive-curvature-maximizer-steady-soliton}
also follows from the linear trace Harnack calculation, with less algebra than
the null-space argument presented here.
\end{remark}

\begin{lemma}[Propagation of a deficient ancient Harnack null space]
\label{lem:chow-ii-ancient-harnack-nullspace-propagation}
\source{chow-et-al-ricci-flow-part-ii:page-0322,chow-et-al-ricci-flow-part-ii:page-0323}
\uses{cor:chow-ii-ancient-matrix-harnack, lem:chow-ii-harnack-exhaustion-function, lem:chow-ii-harnack-barrier-functions}
\dcref{ch15:15.37}
\group{Ancient Solutions and Solitons}


If at one time the null space of \(\overline Q\) has dimension less than
\(n\) at some point, then at every point at every later time its null space
has dimension less than \(n\).
\end{lemma}
\begin{proof}
\sketch On a neighborhood where the nullity is constant, choose a nonzero compactly
supported one-form \(Y\) orthogonal to all one-form projections of null
vectors.  After scaling, \(F_0=Y\otimes Y\) satisfies
\(\overline Q(U\oplus W)\ge F_0(W,W)\).  Evolve \(F_0\) by the tensor heat
equation.  The tensor maximum principle preserves \(F\ge0\), while the strong
maximum principle makes \(\operatorname{tr}F>0\) at all later points.

Apply the noncompact barrier argument used for the matrix Harnack theorem to
\(\overline Q-F\); it remains nonnegative.  Since \(\Rm>0\), a null vector
of \(\overline Q\) is uniquely determined by its one-form component.  Those
components lie in \(\ker F\), whose dimension is at most \(n-1\), proving the
claim.
\end{proof}

\begin{lemma}[The ancient null-space map is wedging with a vector]
\label{lem:chow-ii-nullspace-map-is-wedge-map}
\source{chow-et-al-ricci-flow-part-ii:page-0324,chow-et-al-ricci-flow-part-ii:page-0325,chow-et-al-ricci-flow-part-ii:page-0326}
\uses{lem:chow-ii-ancient-harnack-nullspace-propagation}
\dcref{ch15:15.38}
\group{Ancient Solutions and Solitons}


Suppose the null space of \(\overline Q\) has dimension \(n\), so it is the
graph of \(L:T\mathcal M\to\Lambda^2T\mathcal M\).  Define
\(V_i=2g^{jk}L_{ijk}/(n-1)\).  Then
\[
L_{ijk}=\frac12(V_ig_{jk}-V_jg_{ik}),
\]
so \(L(W)=V\wedge W\).
\end{lemma}
\begin{proof}
\sketch Diagonalize the nonnegative quadratic \(\overline Q\) as
\(\sum_\alpha(X_i^\alpha W^i+Y_{ij}^\alpha U^{ij})^2\).  Its reaction
quadratic is also a sum of squares and vanishes on every null vector.
Substitute \(U=L(W)\) into both families of squares.  After choosing the
\(Y^\alpha\) orthonormally on two-forms, polarization gives
\[
g_{ip}L_{rsq}-g_{iq}L_{rsp}-g_{ir}L_{pqs}+g_{is}L_{pqr}
+g_{qs}L_{pri}-g_{qr}L_{psi}-g_{ps}L_{qri}+g_{pr}L_{qsi}=0.
\]
Contracting \(i,p\) and using the definition of \(V\) yields the formula.
\end{proof}

\begin{remark}[Space-time splitting interpretation]
\label{rem:chow-ii-ancient-rigidity-spacetime-splitting}
\source{chow-et-al-ricci-flow-part-ii:page-0327}
\dcref{ch15:15.39}
\group{Ancient Solutions and Solitons}

The steady soliton equation makes \(\nabla f+\partial_t\) parallel for the
space-time connection.  The ancient Harnack estimate gives nonnegative
space-time curvature, while attainment of the scalar-curvature maximum gives
a null direction.  Theorem
\ref{thm:chow-ii-ancient-positive-curvature-maximizer-steady-soliton} is thus a
space-time splitting theorem.
\end{remark}

\section{Applications and related estimates}

\begin{theorem}[Pinched noncompact three-manifolds are flat]
\label{thm:chow-ii-pinched-noncompact-three-manifold-flat}
\source{chow-et-al-ricci-flow-part-ii:page-0327,chow-et-al-ricci-flow-part-ii:page-0328,chow-et-al-ricci-flow-part-ii:page-0329,chow-et-al-ricci-flow-part-ii:page-0330}
\uses{thm:chow-ii-ancient-positive-curvature-maximizer-steady-soliton, thm:chow-ii-odd-dimensional-positive-curvature-gap}
\dcref{ch15:15.40}
\group{Applications of Harnack Estimates}


Let \((\mathcal M^3,g_0)\) be complete and noncompact, with bounded
nonnegative sectional curvature.  If
\(\Ric_{g_0}\ge\varepsilon R_{g_0}g_0\) for some \(\varepsilon>0\), then
\(g_0\) is flat.
\end{theorem}
\begin{proof}
\sketch Run the complete bounded-curvature Ricci flow from \(g_0\).  The pinching is
preserved and improves to
\(|\Ric-Rg/3|\le CR^{1-\delta}\).  If the initial metric were nonflat, the
strong maximum principle would give positive curvature for positive time.
Finite-time blow-up and point-picking would produce a complete Einstein
three-dimensional blow-up limit; it is a spherical space form, contradicting
noncompactness.  Thus the flow is immortal.

If it is Type IIb, point-picking gives a nonflat eternal limit attaining its
maximum scalar curvature.  Theorem
\ref{thm:chow-ii-ancient-positive-curvature-maximizer-steady-soliton} makes it
a steady gradient soliton, incompatible with the uniform Ricci pinching in
dimension greater than two.  If it is Type III, blow-down gives a nonflat
expanding gradient soliton with positive curvature and exponentially decaying
curvature.  This has zero asymptotic scalar-curvature ratio, contradicting
Theorem~\ref{thm:chow-ii-odd-dimensional-positive-curvature-gap}.  Both cases
are impossible, so \(g_0\) is flat.
\end{proof}

\begin{theorem}[Odd-dimensional positive-curvature gap theorem]
\label{thm:chow-ii-odd-dimensional-positive-curvature-gap}
\source{chow-et-al-ricci-flow-part-ii:page-0330}
\dcref{ch15:15.41}
\group{Applications of Harnack Estimates}

If \((\mathcal M^{2k+1},g)\) is complete, noncompact, and has positive
sectional curvature, then its asymptotic scalar-curvature ratio is positive.
\end{theorem}

\begin{theorem}[Differential Harnack estimate for mean curvature flow]
\label{thm:chow-ii-mean-curvature-flow-harnack}
\source{chow-et-al-ricci-flow-part-ii:page-0330,chow-et-al-ricci-flow-part-ii:page-0331}
\dcref{ch15:15.42}
\group{Related Harnack Estimates}

Let a closed manifold evolve by a strictly convex mean curvature flow
\(\partial_tX=-H\nu\) in Euclidean space.  Then
\[
\partial_tH+\frac{H}{2t}-h^{-1}(\nabla H,\nabla H)\ge0.
\]
\end{theorem}
\begin{proof}
\sketch Denote the left side by \(Z\).  Direct differentiation of the mean-curvature
and second-fundamental-form evolution equations, followed by completing the
quadratic terms, gives
\[
(\partial_t-\Delta)Z
=2g^{ij}(h^{-1})^{k\ell}J_{ik}J_{j\ell}
+\left(|h|^2-\frac2t\right)Z,
\]
where
\[
J_{ik}=\nabla_i\nabla_kH+H(h^2)_{ik}
-(h^{-1})^{pq}\nabla_pH\nabla_qh_{ik}+\frac1{2t}h_{ik}.
\]
Strict convexity makes the first term nonnegative.  The scalar maximum
principle, beginning at small positive time where the \(H/(2t)\) term
dominates, yields \(Z\ge0\).
\end{proof}

\begin{example}[Hamilton's three-dimensional Harnack problem]
\label{ex:chow-ii-arbitrary-three-dimensional-matrix-harnack-problem}
\dcref{ch15:15.32}
\group{Open Harnack Problems}
\source{chow-et-al-ricci-flow-part-ii:page-0320}
Determine whether there is a matrix Harnack estimate for every complete
three-dimensional Ricci flow with bounded curvature, without assuming a
nonnegative curvature operator.
\end{example}

\begin{example}[Complete noncompact mean-curvature-flow problem]
\label{ex:chow-ii-complete-noncompact-mcf-harnack-problem}
\dcref{ch15:15.43}
\group{Open Harnack Problems}
\source{chow-et-al-ricci-flow-part-ii:page-0331}
Find sufficient hypotheses, beyond convexity and bounded second fundamental
form if necessary, under which
Theorem~\ref{thm:chow-ii-mean-curvature-flow-harnack} extends to complete
noncompact mean curvature flows.
\end{example}
